\chapter{Garis, Lingkaran dan Sudut}\label{ch:g6:lines}

Geometri bermula dari penggaris, penggaris siku dan jangka. Bab ini
menetapkan kosakatanya --- titik, \omterm{def:g6:lines:objects}{ruas garis}, sinar garis, garis dan
lingkaran --- memperkenalkan dua kedudukan istimewa garis (\omterm{def:g6:lines:perp}{sejajar} dan
\omterm{def:g6:lines:perp}{tegak lurus}), lalu mengajarkan cara mengukur dan menggambar sudut.

\section{Titik, ruas garis, garis}

\begin{definition}[Objek dasar]\label{def:g6:lines:objects}
Melalui dua titik berbeda $A$ dan $B$ lewat:
\begin{itemize}
  \item \emph{ruas garis}\index{ruas garis} $[AB]$: bagian garis di
        antara $A$ dan $B$ (ia punya panjang, ditulis $AB$);
  \item \emph{sinar garis} $[AB)$: bermula di $A$, melalui $B$ dan
        berlanjut selamanya;
  \item \emph{garis}\index{garis} $(AB)$: memanjang selamanya ke kedua
        arah. Dua titik menentukan tepat satu garis.
\end{itemize}
Titik yang terletak pada garis yang sama disebut \emph{segaris}.
\end{definition}

\begin{omfigure}
\begin{tikzpicture}[scale=1.0]
  \fill (0,1.8) circle (1.5pt) node[above, font=\small] {$A$};
  \fill (2.5,1.8) circle (1.5pt) node[above, font=\small] {$B$};
  \draw[thick, omDef] (0,1.8) -- (2.5,1.8);
  \node[right, font=\small] at (2.9,1.8) {ruas garis $[AB]$};
  \fill (0,0.9) circle (1.5pt) node[above, font=\small] {$A$};
  \fill (2.5,0.9) circle (1.5pt) node[above, font=\small] {$B$};
  \draw[thick, omThm] (0,0.9) -- (4.1,0.9);
  \node[right, font=\small] at (4.3,0.9) {sinar garis $[AB)$};
  \fill (0,0) circle (1.5pt) node[above, font=\small] {$A$};
  \fill (2.5,0) circle (1.5pt) node[above, font=\small] {$B$};
  \draw[thick, omMeth] (-1.4,0) -- (4.1,0);
  \node[right, font=\small] at (4.3,0) {garis $(AB)$};
\end{tikzpicture}

{\small Dua titik yang sama, tiga objek yang berbeda. Kurungnya
mengatakan mana yang berhenti dan mana yang berlanjut: $[$ berhenti, $($ berlanjut.}
\end{omfigure}

\section{Garis sejajar dan garis tegak lurus}

\begin{definition}[Tegak lurus, sejajar]\label{def:g6:lines:perp}
Dua garis disebut \emph{tegak lurus}\index{garis tegak lurus} jika
keduanya berpotongan membentuk \omterm{def:g3:shapes:rightangle}{sudut siku-siku}; kita menulis
$d_1 \perp d_2$ dan menandai \omterm{def:g3:shapes:rightangle}{sudut siku-sikunya} dengan persegi kecil.
Dua garis disebut \emph{sejajar}\index{garis sejajar} jika keduanya
tidak pernah bertemu, sejauh apa pun diperpanjang; kita menulis
$d_1 \parallel d_2$.
\end{definition}

\begin{omfigure}
\begin{tikzpicture}[scale=1.0]
  \draw[thick] (-1.8,0) -- (1.8,0) node[right, font=\small] {$d_1$};
  \draw[thick, omDef] (0.6,-1.2) -- (0.6,1.4)
    node[above, font=\small] {$d_2$};
  \draw[thin] (0.6,0) ++(-0.25,0) -- ++(0,0.25) -- ++(0.25,0);
  \node[below, font=\small] at (0,-1.5) {$d_1 \perp d_2$};
  \begin{scope}[xshift=6.5cm]
  \draw[thick] (-1.8,-0.5) -- (1.8,0.3) node[right, font=\small] {$d_3$};
  \draw[thick, omThm] (-1.8,-1.3) -- (1.8,-0.5)
    node[right, font=\small] {$d_4$};
  \node[below, font=\small] at (0,-1.5) {$d_3 \parallel d_4$};
  \end{scope}
\end{tikzpicture}

{\small Sebuah perpotongan siku-siku (ditandai persegi kecil) dan dua
\omterm{def:g6:lines:perp}{garis sejajar}: arahnya sama, tanpa titik potong.}
\end{omfigure}

\begin{proposition}[Dua fakta yang berguna]\label{prop:g6:lines:facts}
\begin{enumerate}
  \item Jika dua garis sama-sama \omterm{def:g6:lines:perp}{tegak lurus} terhadap garis ketiga,
        keduanya \omterm{def:g6:lines:perp}{sejajar} satu sama lain.
  \item Jika dua \omterm{def:g6:lines:perp}{garis sejajar}, setiap garis yang \omterm{def:g6:lines:perp}{tegak lurus} terhadap
        salah satunya juga \omterm{def:g6:lines:perp}{tegak lurus} terhadap yang lain.
\end{enumerate}
\end{proposition}
\admitted

\begin{method}[Menggambar dengan penggaris siku]\label{met:g6:lines:setsquare}
Untuk menggambar \omterm{def:g6:lines:perp}{garis tegak lurus} terhadap garis $d$ melalui titik $P$:
\begin{enumerate}
  \item letakkan satu tepi \omterm{def:g3:shapes:rightangle}{sudut siku-siku} itu sepanjang garis $d$;
  \item geser penggaris siku itu sepanjang $d$ sampai tepinya yang
        lain mencapai $P$;
  \item gambarlah garis sepanjang tepi itu, lalu tandai \omterm{def:g3:shapes:rightangle}{sudut siku-sikunya}.
\end{enumerate}
Untuk \omterm{def:g6:lines:perp}{garis sejajar} melalui $P$: gambarlah \omterm{def:g6:lines:perp}{garis tegak lurus} terhadap
$d$, lalu \omterm{def:g6:lines:perp}{garis tegak lurus} terhadap \emph{garis itu} melalui $P$
(\cref{prop:g6:lines:facts}).
\end{method}

\section{Lingkaran}

\begin{definition}[Lingkaran]\label{def:g6:lines:circle}
\emph{Lingkaran}\index{lingkaran} berpusat $O$ dengan \omterm{def:g3:shapes:shapes}{jari-jari} $r$
adalah himpunan semua titik yang berjarak tepat $r$ dari $O$. \omterm{def:g6:lines:objects}{Ruas
garis} dari pusat ke lingkaran disebut \emph{\omterm{def:g3:shapes:shapes}{jari-jari}}; \omterm{def:g6:lines:objects}{ruas garis} yang
menghubungkan dua titik lingkaran melalui pusatnya disebut
\emph{\omterm{def:g4:geometry:circle}{diameter}} --- panjangnya $2r$; \omterm{def:g6:lines:objects}{ruas garis} yang menghubungkan dua
titik lingkaran disebut \emph{tali busur}.
\end{definition}

\begin{omfigure}
\begin{tikzpicture}[scale=1.15]
  \draw[thick] (0,0) circle (1.5);
  \fill (0,0) circle (1.5pt) node[below left, font=\small] {$O$};
  \fill (1.5,0) circle (1.5pt) node[right, font=\small] {$M$};
  \draw[omDef, thick] (0,0) -- (1.5,0)
    node[midway, above, font=\small] {jari-jari};
  \fill (-1.06,1.06) circle (1.5pt) node[above left, font=\small] {$A$};
  \fill (1.06,-1.06) circle (1.5pt) node[below right, font=\small] {$B$};
  \draw[omThm, thick] (-1.06,1.06) -- (1.06,-1.06)
    node[pos=0.25, right=2pt, font=\small] {diameter};
  \fill (0.52,1.41) circle (1.5pt) node[above, font=\small] {$C$};
  \fill (1.41,0.52) circle (1.5pt) node[right, font=\small] {$D$};
  \draw[omMeth, thick] (0.52,1.41) -- (1.41,0.52)
    node[midway, above right=-2pt, font=\small] {tali busur};
\end{tikzpicture}

{\small Sebuah lingkaran dengan \omterm{def:g3:shapes:shapes}{jari-jari} $[OM]$, \omterm{def:g4:geometry:circle}{diameter} $[AB]$
(tali busur yang melalui pusat, dua kali panjang \omterm{def:g3:shapes:shapes}{jari-jari}) dan tali busur $[CD]$.}
\end{omfigure}

\begin{example}\label{ex:g6:lines:circle}
``Gambarlah lingkaran berpusat $O$ yang melalui $A$'': bukalah jangka
dari $O$ ke $A$ --- \omterm{def:g3:shapes:shapes}{jari-jarinya} adalah jarak $OA$ --- lalu putar.
Setiap titik lingkaran ini berjarak sama dari $O$ seperti $A$.
\end{example}

\section{Sudut}

\begin{definition}[Sudut]\label{def:g6:lines:angle}
Dua sinar garis $[AB)$ dan $[AC)$ dengan titik pangkal yang sama $A$
membentuk \emph{sudut}\index{sudut} $\widehat{BAC}$; titik $A$ adalah
\emph{\omterm{def:g5:solids:def}{titik sudutnya}} (selalu huruf yang di tengah!). Sudut diukur
dalam \emph{derajat} ($^\circ$), dari $0^\circ$ sampai $360^\circ$
untuk satu putaran penuh. Sebuah sudut disebut:
\begin{itemize}
  \item \emph{siku-siku} jika besarnya $90^\circ$;
  \item \emph{lancip} jika besarnya kurang daripada $90^\circ$;
  \item \emph{tumpul} jika besarnya antara $90^\circ$ dan
        $180^\circ$;
  \item \emph{lurus} jika besarnya $180^\circ$ (kedua sinar garisnya
        membentuk sebuah garis).
\end{itemize}
\end{definition}

\begin{omfigure}
\begin{tikzpicture}[scale=0.9]
  \draw[thick] (0,0) -- (1.8,0);
  \draw[thick] (0,0) -- (40:1.8);
  \draw[omDef, ->] (0.7,0) arc (0:40:0.7);
  \node[omDef, font=\small] at (20:1.0) {$40^\circ$};
  \node[below, font=\small] at (0.9,-0.3) {lancip};
  \begin{scope}[xshift=3.6cm]
  \draw[thick] (0,0) -- (1.8,0);
  \draw[thick] (0,0) -- (0,1.8);
  \draw[thin] (0.3,0) -- (0.3,0.3) -- (0,0.3);
  \node[below, font=\small] at (0.9,-0.3) {siku-siku};
  \end{scope}
  \begin{scope}[xshift=7.2cm]
  \draw[thick] (0,0) -- (1.8,0);
  \draw[thick] (0,0) -- (135:1.8);
  \draw[omThm, ->] (0.6,0) arc (0:135:0.6);
  \node[omThm, font=\small] at (65:1.0) {$135^\circ$};
  \node[below, font=\small] at (0.9,-0.3) {tumpul};
  \end{scope}
  \begin{scope}[xshift=11.6cm]
  \draw[thick] (-1.5,0) -- (1.5,0);
  \fill (0,0) circle (1.2pt);
  \draw[omMeth, ->] (0.5,0) arc (0:180:0.5);
  \node[below, font=\small] at (0,-0.3) {lurus};
  \end{scope}
\end{tikzpicture}

{\small Empat keluarga sudut. \omterm{def:g3:shapes:rightangle}{Sudut siku-siku} ($90^\circ$) adalah
patokannya: lancip berarti lebih kecil, tumpul berarti lebih besar.}
\end{omfigure}

\begin{method}[Mengukur sudut dengan busur derajat]\label{met:g6:lines:protractor}
\begin{enumerate}
  \item Letakkan pusat busur derajat tepat pada \omterm{def:g5:solids:def}{titik sudutnya};
  \item luruskan garis $0^\circ$-nya dengan salah satu sisi sudut itu;
  \item bacalah skala yang dilewati sisi yang lain --- dengan memakai
        skala yang bermula dari $0$ pada sisi yang diluruskan tadi;
  \item periksalah dengan mata: sudut lancip harus terbaca kurang
        daripada $90^\circ$, sudut tumpul lebih.
\end{enumerate}
\end{method}

\begin{example}\label{ex:g6:lines:estimate}
Sebelum mengukur, taksirlah dulu! Sudut yang sedikit lebih terbuka
daripada pojok selembar kertas berukuran sedikit di atas $90^\circ$;
setengah \omterm{def:g3:shapes:rightangle}{sudut siku-siku} adalah $45^\circ$; sepertiga \omterm{def:g3:shapes:rightangle}{sudut siku-siku}
adalah $30^\circ$. Kalau busur derajatmu menyebut $150^\circ$ untuk
sudut yang terlihat lancip, kamu membaca skala yang salah: ukuran yang
benar adalah $180^\circ - 150^\circ = 30^\circ$.
\end{example}

\section{Latihan}

\begin{exercise}[$\star$]\label{exo:g6:lines:1}
Gambarlah tiga titik $A$, $B$, $C$ yang tidak segaris. Gambarlah
dengan warna berbeda: \omterm{def:g6:lines:objects}{ruas garis} $[AB]$, sinar garis $[CA)$, garis $(BC)$.
\end{exercise}

\begin{exercise}[$\star$]\label{exo:g6:lines:2}
Benar atau salah? ``$[AB]$ dan $[BA]$ adalah \omterm{def:g6:lines:objects}{ruas garis} yang sama.''
``$[AB)$ dan $[BA)$ adalah sinar garis yang sama.'' ``$(AB)$ dan
$(BA)$ adalah garis yang sama.'' Jelaskan setiap jawabanmu.
\end{exercise}

\begin{exercise}[$\star$]\label{exo:g6:lines:3}
Gambarlah sebuah garis $d$ dan titik $P$ yang tidak pada $d$.
Lukislah dengan penggaris siku: \omterm{def:g6:lines:perp}{garis tegak lurus} terhadap $d$ melalui
$P$, lalu \omterm{def:g6:lines:perp}{garis sejajar} dengan $d$ melalui $P$. Jelaskan langkahmu.
\end{exercise}

\begin{exercise}[$\star$]\label{exo:g6:lines:4}
Garis $a$ dan $b$ sama-sama \omterm{def:g6:lines:perp}{tegak lurus} terhadap garis $c$, dan garis
keempat $e$ \omterm{def:g6:lines:perp}{tegak lurus} terhadap $a$. Apa yang dapat kamu katakan
tentang $a$ dan $b$? Tentang $e$ dan $c$? Berilah alasan dengan
\cref{prop:g6:lines:facts}.
\end{exercise}

\begin{exercise}[$\star$]\label{exo:g6:lines:5}
Gambarlah lingkaran berpusat $O$ dengan \omterm{def:g3:shapes:shapes}{jari-jari} $3$ cm. Letakkan
titik $M$ pada lingkaran, titik $N$ di dalamnya, titik $P$ di luarnya.
Apa yang dapat kamu katakan tentang jarak $OM$, $ON$, $OP$
dibandingkan dengan $3$ cm?
\end{exercise}

\begin{exercise}[$\star$]\label{exo:g6:lines:6}
Sebuah lingkaran berdiameter $9$ cm. Berapa \omterm{def:g3:shapes:shapes}{jari-jarinya}? Lingkaran
lain berjari-jari $2.6$ cm: berapa \omterm{def:g4:geometry:circle}{diameternya}?
\end{exercise}

\begin{exercise}[$\star$]\label{exo:g6:lines:7}
Sebutkan nama sudut yang ditandai dengan tiga huruf, lalu golongkan
(lancip, siku-siku, tumpul, lurus): sudut $72^\circ$ dengan \omterm{def:g5:solids:def}{titik
sudut} $R$ di antara sinar garis menuju $S$ dan $T$; sudut $148^\circ$
dengan \omterm{def:g5:solids:def}{titik sudut} $B$ di antara sinar garis menuju $A$ dan $C$; sudut
$90^\circ$ dengan \omterm{def:g5:solids:def}{titik sudut} $O$ di antara sinar garis menuju $M$ dan $N$.
\end{exercise}

\begin{exercise}[$\star$]\label{exo:g6:lines:8}
Taksirlah, lalu ukurlah dengan busur derajat, ketiga sudut sebuah
segitiga yang kamu gambar sendiri. Jumlahkan ketiga ukurannya: apa
yang kamu peroleh? (Simpanlah jawabanmu untuk \cref{ch:g7:triangles}.)
\end{exercise}

\begin{exercise}[$\star$]\label{exo:g6:lines:9}
Gambarlah sudut $\widehat{xOy}$ sebesar $65^\circ$ dengan busur
derajat, lalu sudut $130^\circ$. Bagaimana kamu dapat memperoleh sudut
yang kedua dari yang pertama tanpa busur derajat?
\end{exercise}

\begin{exercise}[$\star\star$]\label{exo:g6:lines:10}
Dua desa $A$ dan $B$ digambar pada sebuah peta. Di mana titik yang
berjarak $3$ cm dari $A$ \emph{dan} $2$ cm dari $B$? Gambarlah yang
memperlihatkan ada berapa titik semacam itu ($0$, $1$ atau $2$,
bergantung pada jarak $AB$).
\end{exercise}

\begin{exercise}[$\star\star$]\label{exo:g6:lines:11}
Sebuah jam menunjukkan pukul 3: berapa sudut di antara kedua jarumnya?
Pertanyaan yang sama pada pukul 5, dan pada pukul 6. (Satu putaran
penuh adalah $360^\circ$ untuk $12$ jam.)
\end{exercise}

\begin{exercise}[$\star\star\star$]\label{exo:g6:lines:12}
Gambarlah \omterm{def:g6:lines:objects}{ruas garis} $[AB]$ sepanjang $6$ cm. Lukislah titik $C$
sedemikian sehingga $AC = BC = 6$ cm, hanya dengan jangka, lalu ukurlah
sudut $\widehat{CAB}$. Segitiga apa yang kamu bangun, dan apa dugaanmu
tentang sudutnya?
\end{exercise}

\section{Soal: Geometri pada muka jam}

\begin{problem}[{Soal akhir pekan --- sudut sebagai pecahan dari satu
putaran: membacanya pada jam, memburu saat kedua jarum bertemu, lalu
mengiris satu hari menjadi diagram lingkaran}]
\label{pb:g6:lines:1}
Sebuah jam adalah busur derajat yang menunjukkan waktu: mukanya
adalah satu putaran penuh $360^\circ$, yang dipotong oleh dua belas
tanda jam menjadi dua belas sudut yang sama. \Cref{exo:g6:lines:11}
mengukur jarumnya pada pukul 3 dan pukul 6; soal ini membangun
teorinya secara lengkap --- termasuk saat ketika \emph{tidak satu
pun} jarum menunjuk sebuah tanda --- menjawab pertanyaan yang jarang
dijawab benar oleh orang dewasa (``berapa sering kedua jarum duduk
tepat bertumpuk?''), lalu ditutup dengan mengiris satu hari penuh
menjadi diagram lingkaran.

\medskip
\noindent\textbf{Bagian I --- \omterm{def:g6:fractions:def}{Pecahan} dari satu putaran.}
\begin{enumerate}
  \item Satu putaran penuh berukuran $360^\circ$. Berapa derajat
        setengah putaran, \omterm{def:g3:division:half}{seperempat} putaran, dan seperdua belas
        putaran?
  \item Kedua belas tanda jam memotong muka jam itu menjadi dua
        belas sudut yang sama besar di pusatnya. Berapa derajat di antara dua tanda yang
        bertetangga? Perolehlah kembali jawaban
        \cref{exo:g6:lines:11} untuk pukul 3 dengan membilang tandanya.
  \item Sebutkan sudut di antara kedua jarum pada pukul 1, pukul 4 dan
        pukul 7. (Pada pukul 7 kedua jarum memisahkan mukanya menjadi
        dua sudut; ``sudut di antara kedua jarum'' selalu berarti yang
        \emph{lebih kecil}.)
  \item Jarum menit itu menempuh satu putaran penuh dalam $60$ menit.
        Berapa derajat yang disapunya setiap menit?
  \item Jarum jam berjalan dari satu tanda ke tanda berikutnya ---
        $30^\circ$ --- dalam $60$ menit. Berapa derajat yang disapu
        \emph{jarum itu} setiap menit? (Sebuah \omterm{def:g6:decimals:places}{bilangan desimal},
        \cref{ch:g6:decimals}.)
\end{enumerate}

\medskip
\noindent\textbf{Bagian II --- Saat ketika tidak ada yang menunjuk
sebuah tanda.}
\begin{enumerate}[resume]
  \item Pada pukul 3:30, jarum menit menunjuk angka $6$. Di mana
        tepatnya jarum jam? Hitunglah sudut setiap jarum dari angka
        $12$ (diukur searah jarum jam), lalu simpulkan sudut di antara
        kedua jarum.
  \item Pada pukul 6:30 banyak orang menduga kedua jarumnya bertumpuk.
        Terkalah dulu, lalu hitunglah sudutnya seperti pada pertanyaan
        6. Siapa yang benar?
  \item Hitunglah sudut di antara kedua jarum pada pukul 9:15.
  \item Jelaskan mengapa, di suatu saat antara pukul 1:00 dan 1:10,
        kedua jarum pasti bertumpuk tepat: di mana jarum menit
        terhadap jarum jam pada pukul 1:00, dan di mana pada pukul
        1:10? Jarum mana yang berjalan lebih cepat, dan mengapa hal
        itu menyelesaikannya?
  \item Dalam dua belas jam, berapa kali kedua jarum duduk tepat
        bertumpuk? (Satu pertemuan terjadi pada setiap rentang antara
        jam yang berurutan --- dengan satu kekecualian: apa yang
        terjadi antara pukul 11 dan 12? Sebutkan waktu pertemuan yang
        kira-kira lalu bilanglah.) Berapa kali dalam satu hari penuh?
\end{enumerate}

\medskip
\noindent\textbf{Bagian III --- Sudut di sekeliling sebuah titik:
diagram lingkaran.}
\begin{enumerate}[resume]
  \item Kedua belas juring pada muka jam bersama-sama mengisi seluruh
        mukanya: sudutnya berjumlah $360^\circ$. Jelaskan mengapa hal
        ini benar untuk kumpulan sudut \emph{mana pun} yang berbagi
        satu \omterm{def:g5:solids:def}{titik sudut} dan mengisi satu putaran penuh di
        sekelilingnya, tanpa tumpang tindih dan tanpa celah.
  \item Zahra mencatat harinya yang $24$ jam: tidur $9$ jam, sekolah
        $6$ jam, bermain $3$ jam, makan $2$ jam, selebihnya $4$ jam.
        Ia menginginkan sebuah \emph{diagram lingkaran}: sebuah cakram
        yang setiap kegiatannya mendapat satu juring, dengan sudut
        yang sebanding dengan waktunya. Berapa \omterm{def:g6:fractions:def}{pecahan} hari untuk
        setiap kegiatan, dan berapa derajat untuk juringnya? Periksa
        bahwa kelima sudutnya berjumlah $360^\circ$.
  \item Jelaskan, selangkah demi selangkah, cara menggambar diagram
        lingkaran Zahra dengan jangka dan busur derajat
        (\cref{met:g6:lines:protractor}): ke mana \omterm{def:g3:shapes:shapes}{jari-jari} yang
        pertama diarahkan, dan bagaimana setiap juring baru bermula di
        tempat juring sebelumnya berakhir.
  \item Pada diagram lingkaran sebuah hari yang lain, satu juringnya
        berukuran tepat $90^\circ$. Berapa \omterm{def:g6:fractions:def}{pecahan} cakramnya, dan
        berapa jam yang diwakilinya?
  \item Penutupnya, kembali ke jam: dalam $24$ jam, berapa putaran
        penuh yang ditempuh jarum jam? Jarum menit? Jarum detik?
        (Salah satu jawabannya lebih dari seribu.)

\end{enumerate}
\end{problem}
